% Summer research report: two-particle Hamiltonian system
% Restructured 2026-11-18 around the regular fraction f(E). Earlier versions:
%   report/archive/report_2026-11-07_before_simplify.tex
%   report/archive/report_2026-11-18_before_restructure.tex
% Data sources: report01-04.ipynb, case3_newton_failure_diagnosis.ipynb,
% lyapunov_*.py, regular_fraction_scan.py; figures from the scripts in Table 2.
% Text marked with \todo{...} still needs checking or a decision.

\documentclass[11pt,a4paper]{article}

\usepackage[margin=2.5cm]{geometry}
\usepackage{amsmath,amssymb}
\usepackage{graphicx}
\usepackage{booktabs}
\usepackage{xcolor}
\usepackage[colorlinks=true,linkcolor=blue!50!black,citecolor=blue!50!black,urlcolor=blue!50!black]{hyperref}

\graphicspath{{../images/}{figures/}}

% Figures and tables use [H] (float package): placed exactly where they are
% in the source, right after the paragraph that first refers to them.
\usepackage{float}

% Visible placeholder. For a clean draft use \newcommand{\todo}[1]{}
\newcommand{\todo}[1]{\textcolor{red}{[TODO: #1]}}

\newcommand{\q}{\mathbf{q}}
\newcommand{\p}{\mathbf{p}}
\newcommand{\Vs}{V_{\mathrm{sym}}}
\newcommand{\Eeq}{E_{\mathrm{eq}}}

\title{Regular and Chaotic Motion of Two Repelling Particles in a Box:\\
a Numerical Study}
\author{Jiajun Wu\\[2pt]
\small Department of Mathematics, Imperial College London\\
\small Summer research project (2026), supervised by Prof.\ Dmitry Turaev}
\date{November 2026}

\begin{document}
\maketitle

% ---------------------------------------------------------------------------
\begin{abstract}
We study numerically two particles in a square box that repel each other and
are pushed back by the walls, restricting to motions that are symmetric about
the centre of the box. Our main result is the fraction $f(E)$ of the energy
surface filled by regular motion, measured with finite-time Lyapunov exponents
for initial conditions sampled uniformly on the energy surface. Near the
bottom of the four potential wells all motion is regular. As the energy
approaches $E_s\approx1.9969$, the energy of the saddles between the wells,
chaos appears, and around $E_s$ the regular fraction drops sharply, to about
$4\%$ at $E\approx2.05$; it then rises slowly, to about one third at $E=200$. The regular motion
at high energy lies on islands around stable periodic orbits. We follow two of
these orbits in energy with Newton's method. A ``vertical'' orbit, of the type
studied by Rom-Kedar and Turaev, is linearly stable for $E\gtrsim3.9$ and
unstable below. A ``circular'' orbit around the centre is linearly stable at
every energy we computed, from $E\approx2$ to $E\approx202$; at high energy it
approaches a diamond-shaped billiard orbit.
\end{abstract}

% ---------------------------------------------------------------------------
\section{Introduction}\label{sec:intro}

Two particles that repel each other inside a box form a very simple system,
but their long-time motion is not obvious. They might settle into a repeating
pattern, or move chaotically over the whole box, depending on the energy and
the initial condition.

Rom-Kedar and Turaev \cite{turaev2022} proved that, under mild conditions,
the motion of $N$ repelling particles in a bounded region is not ergodic at
sufficiently high energies. They construct stable periodic orbits that are
surrounded by a set of positive measure of quasi-periodic motions; trajectories
that start in this set stay near the periodic orbit for all time, so they
cannot explore the whole energy surface. For particles in a rectangular box,
the stable orbit they construct is one in which every particle moves fast in
the same direction, say vertically, each on its own line and in synchrony with
the others \cite[Sec.~2.6]{turaev2022}. Their result is asymptotic: it holds
only above some energy that is not given explicitly, and it needs the wall
potential to be steep enough (Section~\ref{sec:vertical}). It therefore says
nothing about moderate energies, and nothing about how much of the energy
surface is regular.

In this project we study numerically the simplest case, two particles in a
square box, restricted to motions that are symmetric about the centre. We ask
how the fraction of regular motion on the energy surface depends on the
energy. One might expect regular motion to dominate at both ends: near the
equilibrium, where the motion is a small oscillation, and at high energy,
where the particles move almost like billiard balls in a square, whose
single-particle motion is integrable. We test this by measuring the regular
fraction $f(E)$ from the equilibrium up to $E=200$. We then explain its shape:
at low energy by the shape of the potential, and at high energy by stable
periodic orbits, two families of which we follow in energy.

The report is organised as follows. Section~\ref{sec:model} describes the
model and its potential landscape, and Section~\ref{sec:methods} the numerical
methods. Section~\ref{sec:results} presents the results: examples of the three
kinds of motion, the regular fraction $f(E)$, the transition at low energy,
the phase space at high energy, and the vertical and circular periodic orbits.
Section~\ref{sec:discussion} summarises the findings and their limitations.
All results are numerical; none of them is proved.

% ---------------------------------------------------------------------------
\section{The model}\label{sec:model}

Two particles with positions $\q_i=(x_i,y_i)$ and momenta $\p_i$ ($i=1,2$)
move in the box $(0,L_x)\times(0,L_y)$. The Hamiltonian is
\begin{equation}\label{eq:H}
  H(\q,\p)=\sum_{i=1}^{2}\frac{|\p_i|^2}{2m_i}
  + V\bigl(|\q_1-\q_2|\bigr) + U(\q_1) + U(\q_2),
\end{equation}
where
\begin{equation}\label{eq:potentials}
  V(r)=k\,r^{-\alpha},\qquad
  U(x,y)=C\Bigl[x^{-\beta}+(L_x-x)^{-\beta}+y^{-\beta}+(L_y-y)^{-\beta}\Bigr].
\end{equation}
$V$ is the repulsion between the particles and $U$ pushes each particle away
from the walls. Both become infinite at collision or at a wall, so the
particles can never collide or leave the box.

We take $k=C=\alpha=\beta=1$, $m_1=m_2=1$ and $L_x=L_y=10$. These are simply
convenient values of order one; the model is not meant to describe a specific
physical system.

\paragraph{Energy.} The energy $H$ is conserved. Because of the walls,
momentum and angular momentum are not.

\paragraph{Symmetric motions.}
If particle~2 starts at the point reflection of particle~1 through the centre
of the box, $\q_2=(10,10)-\q_1$, with the opposite momentum $\p_2=-\p_1$, then
by symmetry the two particles stay in this position relative to each other for
all time. We call such motions \emph{centre-symmetric}. For them it is enough
to follow particle~1. Its motion is a Hamiltonian system with two degrees of
freedom, and the total energy is
\begin{equation}\label{eq:Hsym}
  H_{\mathrm{sym}}(\q_1,\p_1)=|\p_1|^2+\Vs(\q_1),\qquad
  \Vs(\q_1)=V\bigl(|2\q_1-(10,10)|\bigr)+2\,U(\q_1),
\end{equation}
(the Hamiltonian that generates the motion of particle~1 is
$\tfrac12H_{\mathrm{sym}}$). All initial conditions in this report are
centre-symmetric.

\paragraph{Wells and saddles.}
Figure~\ref{fig:potential} shows $\Vs$. By the symmetry of the square it has
four equivalent minima, at $\q_1=(a,a)$, $(10-a,a)$, $(a,10-a)$ and
$(10-a,10-a)$ with $a\approx3.6645$, all with energy
$\Eeq\approx1.98765$; this is the equilibrium, with the particles at rest on a
diagonal of the box. Between neighbouring minima there are four saddles, at
$\q_1\approx(3.1926,5)$ and its images, with energy $E_s\approx1.99689$. The
centre of the box, where the particles would collide, is a singularity. So the
bottom of the potential is a ring-shaped valley around the centre, with four
wells separated by saddles only $E_s-\Eeq\approx0.0092$ higher
(Figure~\ref{fig:potential}(d)). For $E<E_s$ the region $\{\Vs<E\}$ that
particle~1 can reach consists of four separate pieces, one around each
minimum; for $E>E_s$ they join into a ring (Figure~\ref{fig:potential}(b,c)).
% Values from images/make_potential_figure.py (2026-11-18): well Hessian
% eigenvalues 0.0228, 0.2454; saddle Hessian eigenvalues -0.0207, 0.3050.

% Generated by images/make_potential_figure.py
\begin{figure}[H]
  \centering
  \includegraphics[width=\textwidth]{fig_potential.pdf}
  \caption{The potential $\Vs(\q_1)$ of a centre-symmetric motion. (a)
  Contours near the bottom; black dots are the four minima, red crosses the
  four saddles, and the white line is the level $\Vs=E_s$. (b), (c) The region
  $\{\Vs<E\}$ that particle~1 can reach just below and just above $E_s$.
  (d) The bottom of the valley, $\min_\rho\Vs$ along each ray from the centre,
  against the angle of the ray. Panels (a)--(c) show only the central part of
  the box.}
  \label{fig:potential}
\end{figure}

% ---------------------------------------------------------------------------
\section{Numerical methods}\label{sec:methods}

\subsection{Time integration}
We integrate the equations of motion with the leapfrog method:
\begin{equation}\label{eq:leapfrog}
  \p^{n+1/2}=\p^n+\tfrac h2 F(\q^n),\quad
  \q^{n+1}=\q^n+h\,\p^{n+1/2},\quad
  \p^{n+1}=\p^{n+1/2}+\tfrac h2 F(\q^{n+1}),
\end{equation}
where $F=-\nabla(V+U)$ is the force and $h$ the time step. Leapfrog is a
second-order method whose energy error stays small and does not grow over long
times \cite{hairer2006}, which is important because we integrate for up to
$10^4$ time units. We use smaller steps when the particles move faster, from
$h=10^{-2}$ near the equilibrium down to $h=10^{-5}$ at $E\approx202$. For the
periodic orbits the relative energy error is at most about $10^{-6}$, and
halving $h$ changes the computed periods by less than one part in $10^{7}$.
For the regular-fraction scan (Section~\ref{sec:fraction-method}) we use
$h=\min\bigl(5\times10^{-3},\,0.026/E\bigr)$, which keeps the relative energy
error below $6\times10^{-4}$.
% Step-halving check run 2026-11-07: E=11.075 T=6.87461211/...275/...291;
% E=3.876 T=14.80128691/...697/...699; E=3.011 T=23.80233268/...379/...407.
% Scan: max|dE|/E = 1.7e-4 (E=10), 1.6e-4 (E=52.23), 6.4e-4 (E=200).

\subsection{Poincar\'e sections}
A Poincar\'e section records the state each time the orbit crosses a chosen
surface in a chosen direction, for example $y_1=5$ with $p_{y_1}>0$. The full
phase space is 8-dimensional; fixing the energy and the surface leaves a
6-dimensional section, which is the one Newton's method below works on. For a
centre-symmetric motion, however, the state of particle~2 is fixed by that of
particle~1, so the motion lies in the 4-dimensional space
$(x_1,y_1,p_{x_1},p_{y_1})$. On a fixed energy level and the surface $y_1=c$,
$p_{y_1}>0$, a point is then determined by $(x_1,p_{x_1})$, because $p_{y_1}$
follows from the energy. The section of a centre-symmetric motion is therefore
exactly two-dimensional, not a projection. In the same way, the surface
$x_1=c$, $p_{x_1}>0$ gives the coordinates $(y_1,p_{y_1})$. Since a time step
usually jumps over the surface, we locate each crossing by linear
interpolation between the two neighbouring steps, or, where Newton's method
needs it more accurately, by a shortened last leapfrog step that lands on the
surface.

On a two-dimensional section the pattern of the points tells us the type of
motion \cite{lichtenberg1992}: a periodic orbit that crosses the section $k$
times per period gives $k$ points, a regular (quasi-periodic) orbit gives
points on one or several smooth closed curves, and a chaotic orbit gives points
scattered over a two-dimensional region. This reading relies on the section
being two-dimensional: in a projection of a higher-dimensional section,
quasi-periodic motion could also look scattered.

\subsection{Lyapunov exponents}\label{sec:lyap-method}
The Lyapunov exponent $\lambda_{\max}$ of a trajectory measures how fast two
nearby trajectories separate: it is $0$ for regular motion and positive for
chaotic motion. We estimate it by following, along the trajectory, a small
perturbation with the linearised leapfrog map, renormalising it every 40 steps
and averaging the logarithm of its growth over a time $T$. The perturbation is
kept centre-symmetric, so $\lambda_{\max}$ measures instability within the
centre-symmetric motions only.

To decide whether a trajectory is regular we need a threshold for
$\lambda_T$. For a regular trajectory nearby trajectories separate only
linearly in time, so the finite-time estimate $\lambda_T$ is not exactly $0$
but of order $\ln T/T$, and it decreases as $T$ grows. For a chaotic one they
separate exponentially, and $\lambda_T$ tends to a positive constant. We
therefore call a trajectory \emph{regular} if
\begin{equation}\label{eq:threshold}
  \lambda_T < c\,\frac{\ln T}{T},\qquad c=3.2 .
\end{equation}

To see whether the result depends on $c$, we looked at the values of
$\lambda_T$, in units of $\ln T/T$, for all trajectories of the
regular-fraction scan of Section~\ref{sec:fraction-method}: $3600$
trajectories at each of $23$ energies from $E=1.99$ to $200$.
At most energies only very few trajectories (at most $2.5\%$) have
$\lambda_T$ between $2\ln T/T$ and $5\ln T/T$, so any $c$ between $2$ and $5$
gives almost the same classification; Figure~\ref{fig:threshold} shows this
gap at three energies. The exception is $E=1.995$--$1.997$, just below and at
$E_s$, where the chaos is weakest, with up to $9\%$ of the trajectories in the
gap. We take
$c=3.2\approx\sqrt{2\cdot5}$, the middle of the gap on the logarithmic scale,
and report the results for $c=2$ and $c=5$ as well.
% Statistics from images/make_threshold_figure.py (standard T, 2026-11-18):
% fraction with 2<u<5: <= 0.025 except E=1.995 (0.078), 1.996 (0.094),
% 1.997 (0.033).

% Generated by images/make_threshold_figure.py
\begin{figure}[H]
  \centering
  \includegraphics[width=\textwidth]{fig_threshold.pdf}
  \caption{Histograms of $\lambda_T$ in units of $\ln T/T$ for the $3600$
  trajectories of the regular-fraction scan at three energies (standard $T$).
  The left group is regular, the right group chaotic. Red line: the threshold
  $c=3.2$; blue band: $2\le c\le5$.}
  \label{fig:threshold}
\end{figure}

\subsection{The regular fraction}\label{sec:fraction-method}
We define the regular fraction $f(E)$ as the fraction of the energy surface
$\{H_{\mathrm{sym}}=E\}$, measured with the Liouville (microcanonical)
measure, that is filled by regular trajectories. To estimate it we sample
initial conditions from this measure, integrate each one and count the regular
ones.

For the two-degree-of-freedom Hamiltonian~\eqref{eq:Hsym} the sampling is
simple. Integrating the microcanonical density over the momenta gives
$\int\delta\bigl(E-|\p_1|^2-\Vs(\q_1)\bigr)\,d^2\p_1=\pi$, which does not depend
on $\q_1$. So we draw $\q_1$ uniformly from the accessible region
$\{\Vs<E\}$, the direction of $\p_1$ uniformly from $[0,2\pi)$, and set
$|\p_1|=\sqrt{E-\Vs(\q_1)}$. We use $3600$ trajectories per energy, so the
statistical error of $f$ is at most about $\pm0.008$.

We do not sample on a Poincar\'e section for two reasons. First, a section
misses every trajectory that never crosses it: no trajectory reaches $y_1=5$
for $E<E_s$, and at high energy the horizontal analogue of the vertical orbit
never does. Second, area on the section is not the same as volume on the
energy surface: the volume filled by a set of trajectories is its area on the
section weighted by the return time.

For $E>2.05$ the integration time is chosen so that every energy gets a
similar number of wall-to-wall flights:
$T=\min\bigl(2500,\;500\sqrt{52.23/E}\bigr)$. For $E\le2.05$, near the
equilibrium and the saddles, the motion is a slow oscillation rather than a
flight between the walls, and the chaos is very weak, so we use the longer
time $T=10\,000$. Section~\ref{sec:fraction} checks that these times are long
enough.

\subsection{Finding periodic orbits}\label{sec:newton}
Let $P$ be the Poincar\'e map, which sends a point $z$ on the section to the
point where the orbit next crosses it. At a fixed energy, $z$ has six
coordinates. A periodic orbit is a fixed point, $P(z^*)=z^*$ (or
$P^k(z^*)=z^*$ if the orbit crosses the section $k$ times per period). We solve
this equation with Newton's method,
\begin{equation}\label{eq:newton}
  z_{j+1}=z_j-\bigl(DP(z_j)-I\bigr)^{-1}\bigl(P(z_j)-z_j\bigr),
\end{equation}
where the Jacobian $DP$ is computed by finite differences. If a full Newton
step does not reduce the error, we take a shorter step. From a good initial
guess Newton converges in 3--6 iterations, with $|P(z)-z|$ between $10^{-9}$
and $10^{-13}$.

\paragraph{Stability.}
The eigenvalues $\lambda$ of $DP(z^*)$, called Floquet multipliers, tell us
what happens to a small perturbation of the periodic orbit after each period:
it is multiplied by roughly $|\lambda|$ in the corresponding direction. If all
multipliers satisfy $|\lambda|=1$, small perturbations neither grow nor shrink
and the orbit is \emph{linearly stable}. If some $|\lambda|>1$, perturbations
grow and the orbit is \emph{unstable}. Since our Jacobians are accurate to
about $10^{-6}$, we count $|\lambda|$ within $10^{-3}$ of $1$ as equal to $1$.

\paragraph{Following an orbit in energy.}
To see how a periodic orbit changes with energy, we change $E$ in small steps
and use the orbit found at the previous energy, with its momentum rescaled, as
the starting guess for Newton at the new energy.

% ---------------------------------------------------------------------------
\section{Results}\label{sec:results}

\subsection{Three kinds of motion}\label{sec:sections}
We first look at single trajectories. We start from the position
$\q_1=(a,a)+0.25\,(1,1)$ near a minimum and give the particles vertical
momenta of increasing size (Figure~\ref{fig:energy-sections}).
\begin{itemize}
  \item At $E\approx2.007$ the section points lie on smooth curves: regular
        motion.
  \item At $E\approx2.028$ and $E\approx2.97$ the points are scattered and the
        trajectory wanders over a large part of the box: chaotic motion.
  \item At $E\approx11$ the motion is regular again. The particles stay in two
        thin vertical strips, close to the stable vertical periodic orbit of
        Section~\ref{sec:vertical}.
\end{itemize}
These are single trajectories, chosen by hand. They show what regular and
chaotic motion look like, but not how common each kind is; that is measured in
the next section.

\begin{figure}[H]
  \centering
  \includegraphics[width=\textwidth]{fig_energy_sections.pdf}
  \caption{Starting from $\q_1=(a,a)+0.25\,(1,1)$ with vertical momenta
  $\pm v$ at four energies: trajectories (top) and 200 returns to the section
  $y_1=a$, $p_{y_1}>0$ (bottom).}
  \label{fig:energy-sections}
\end{figure}

\subsection{The regular fraction}\label{sec:fraction}
Figure~\ref{fig:fraction} shows the regular fraction $f(E)$ at 23 energies
from $E=1.99$ to $E=200$. The curve has three stages.
\begin{itemize}
  \item Near the bottom of the wells all trajectories are regular: $f=1$ at
        $E=1.990$. Already below $E_s$, however, $f$ starts to fall:
        $f\approx0.94$ at $E=1.995$ and $0.67$ at $E=1.996$.
  \item Around and just above $E_s$ the regular fraction drops steeply:
        $f\approx0.45$ at $E=1.997$, $0.30$ at $E=2.000$, $0.14$--$0.21$ for
        $E=2.005$--$2.02$, and a minimum of $f\approx0.04$ at $E=2.05$.
  \item Above that, $f$ rises again, to about $0.20$--$0.27$ for
        $E=2.2$--$3$, and then slowly, to $0.28$ at $E=10$ and about
        $0.32$--$0.34$ for $E=52$--$200$.
\end{itemize}
So the expectation of Section~\ref{sec:intro} holds near the equilibrium, but
only in a very narrow range of energies, and it does not hold, up to $E=200$,
at high energy: there the regular fraction grows only slowly, and about two
thirds of the energy surface is still chaotic.

% Generated by images/make_fraction_figure.py
\begin{figure}[H]
  \centering
  \includegraphics[width=0.8\textwidth]{fig_regular_fraction.pdf}
  \caption{Regular fraction $f(E)$ of the (centre-symmetric) energy surface,
  against $E-\Eeq$ (bottom axis) and $E$ (top axis), with the threshold
  $c=3.2$ of~\eqref{eq:threshold} and the integration times of
  Section~\ref{sec:fraction-method}. Each point uses $3600$ trajectories
  sampled uniformly on the energy surface; the statistical error is at most
  $\pm0.008$. The lines between points only guide the eye. Dashed line: the
  saddle energy $E_s$.}
  \label{fig:fraction}
\end{figure}

\paragraph{Dependence on the threshold.}
For $E\ge1.997$ the thresholds $c=2$ and $c=5$ give values of $f$ within
$0.03$ of those in Figure~\ref{fig:fraction}. At $E=1.995$ and $1.996$, where
the chaos is weakest, the dependence is larger: $f=0.90$--$0.98$ and
$0.64$--$0.73$.

\paragraph{Is $T$ long enough?}
Table~\ref{tab:convergence} compares, for the same trajectories, the regular
fraction with the standard integration time $T$ and with $2T$, at every
energy. A longer time turns some trajectories that looked regular into chaotic
ones; the opposite happened only once (one trajectory at $E=1.995$), so in
practice $f$ can only decrease as $T$ grows. These are trajectories that stay
near an island for a long time before escaping.

For $E>2.05$ the change is at most $0.02$, comparable to the statistical
error, so the curve in Figure~\ref{fig:fraction} is reliable there.

For $E\le2.05$ the time $T=2500$ given by the rule for higher energies is far
too short: at $E=1.996$, for example, it gives $f=0.996$, against $0.667$ with
$T=10\,000$, and at $E=1.997$ it gives $0.631$ against $0.448$. This is why we
use $T=10\,000$ there. Going on from $T=10\,000$ to $20\,000$ changes $f$ much
less: by $0.044$ at $E=1.995$ and $0.037$ at $E=1.996$, by about $0.02$ at
$E=1.997$ and $2.01$, and by at most $0.015$ at the other energies. So the
values for $E\le2.05$ are close to converged; just below $E_s$ they may still
be a few hundredths too high.

\begin{table}[H]
\centering
\caption{Regular fraction $f$ (threshold $c=3.2$) for the same $3600$
trajectories with the standard integration time $T$ and with $T'=2T$.
``Changed'' is the number of trajectories classified differently; all of them
are regular with $T$ and chaotic with $2T$, except one trajectory at
$E=1.995$ (chaotic with $T$, regular with $2T$).}
\label{tab:convergence}
\begin{tabular}{@{}rrrrrr@{}}
\toprule
$E$ & $T$ & $f(T)$ & $T'$ & $f(T')$ & changed\\
\midrule
1.990 & 10\,000 & 1.000 & 20\,000 & 1.000 & 0\\
1.995 & 10\,000 & 0.939 & 20\,000 & 0.895 & 159\\
1.996 & 10\,000 & 0.667 & 20\,000 & 0.630 & 134\\
1.997 & 10\,000 & 0.448 & 20\,000 & 0.428 & 73\\
1.998 & 10\,000 & 0.354 & 20\,000 & 0.339 & 51\\
1.999 & 10\,000 & 0.315 & 20\,000 & 0.301 & 52\\
2.000 & 10\,000 & 0.298 & 20\,000 & 0.292 & 20\\
2.005 & 10\,000 & 0.199 & 20\,000 & 0.189 & 37\\
2.010 & 10\,000 & 0.207 & 20\,000 & 0.188 & 71\\
2.020 & 10\,000 & 0.141 & 20\,000 & 0.134 & 24\\
2.050 & 10\,000 & 0.036 & 20\,000 & 0.036 & 1\\
\midrule
2.10  & 2494 & 0.117 & 4987 & 0.116 & 4\\
2.20  & 2436 & 0.233 & 4872 & 0.229 & 15\\
2.40  & 2333 & 0.223 & 4665 & 0.221 & 7\\
2.70  & 2199 & 0.266 & 4398 & 0.257 & 31\\
3.00  & 2086 & 0.204 & 4173 & 0.204 & 1\\
4.00  & 1807 & 0.247 & 3614 & 0.236 & 43\\
6.00  & 1475 & 0.281 & 2950 & 0.262 & 69\\
10.0  & 1143 & 0.281 & 2285 & 0.277 & 16\\
20.0  & 808  & 0.311 & 1616 & 0.299 & 42\\
52.23 & 500  & 0.328 & 1000 & 0.315 & 44\\
100   & 361  & 0.316 & 723  & 0.306 & 35\\
200   & 256  & 0.337 & 511  & 0.323 & 48\\
\bottomrule
\end{tabular}
\end{table}

None of this changes the shape of $f(E)$: the steep drop around $E_s$, the
minimum at $E=2.05$, the dip at $E=3$ and the slow rise at high energy all
remain with the longer times.

\paragraph{Is $f$ still growing at high energy?}
From $E=10$ to $E\approx52$ the regular fraction clearly grows, from $0.28$ to
$0.33$. Above that it stays at about $0.32$--$0.34$: the values at $E\approx52$,
$100$ and $200$ differ by about $0.01$--$0.02$, which is only one to two times
the statistical error of each value ($\pm0.008$), and they do not increase
steadily ($0.328$, $0.316$, $0.337$). So we cannot tell whether $f$ is still
growing slowly above $E\approx52$ or has levelled off; larger samples and
higher energies would be needed.

\subsection{The transition at low energy}\label{sec:transition}
The steep drop of $f$ happens around the saddle energy $E_s$
(Figure~\ref{fig:potential}). Near the bottom of a well the potential is
nearly quadratic, so the motion is close to two independent harmonic
oscillations, which is an integrable system; the nonlinear terms are small,
and by KAM theory almost all trajectories stay on invariant tori
\cite{lichtenberg1992}. This explains $f=1$ at $E=1.990$.

As $E$ approaches $E_s$, the part of a well that particle~1 can reach extends
up to the saddles (Figure~\ref{fig:potential}(b)). Near a saddle the particle
moves slowly, and what it does next depends very sensitively on how it
arrives: it may come close to passing over the saddle and then turn back.
Motion near the separatrices of a saddle is the classic source of chaos in
Hamiltonian systems \cite{lichtenberg1992}, and the chaotic layer around a
separatrix has a finite width in energy, on both sides of the saddle energy.
This is what we see. At $E=1.995$ and $1.996$ no trajectory leaves its well
(none crosses the line $y_1=5$), yet about $6\%$ and $33\%$ of them are
chaotic. In a sample of nine trajectories at $E=1.996$, the chaotic ones were
those that reached towards the saddles, while most regular ones stayed near
the bottom of the well; we have not checked this systematically.

For $E>E_s$ particle~1 can also pass over the saddles from one well to the
next, and the chaotic regions of the four wells join. Because the barrier is
so low, $E_s-\Eeq\approx0.0092$, most trajectories soon come near a saddle:
the fraction of trajectories that cross $y_1=5$ (that is, pass over a saddle)
grows from $0$ below $E_s$ to $34\%$ at $E=1.997$, $72\%$ at $E=2.000$ and
over $99\%$ at $E=2.05$, while $f$ falls to its minimum, $0.04$ at $E=2.05$.
This explanation is an interpretation of the data, not a proof.

As $E$ grows further, the wells and saddles become irrelevant compared with the
total energy, and new stable periodic orbits appear, such as the vertical orbit
of Section~\ref{sec:vertical}. Their islands provide the regular motion at high
energy, which we look at next.

\subsection{Chaotic region and stable islands at \texorpdfstring{$E\approx52$}{E = 52}}\label{sec:lyapunov}
Figure~\ref{fig:lyapunov} shows the Poincar\'e section $y_1=5$, $p_{y_1}>0$ at
$E\approx52.23$, with every point coloured by the Lyapunov exponent of its
trajectory. For this picture the initial conditions are a uniform grid of
13\,920 points on the section, integrated for $T=500$. Most of the section is
a chaotic region with $\lambda_{\max}\approx0.15$--$0.27$. Inside it there are
many islands of regular motion with $\lambda_{\max}\approx0$. The histogram on
the right has two clear peaks, one for regular and one for chaotic
trajectories.

The islands cover $33\%$ of the section area, or $31\%$ when each trajectory is
weighted by its return time. Sampling on the energy surface gives
$f\approx0.33$ with the same $T=500$ (Section~\ref{sec:fraction}). The three numbers are close, so at
this energy the section is a fair picture of the whole energy surface.

At the centre of each island there is a stable periodic orbit.
Figure~\ref{fig:islands} shows ten of them, found with Newton's method. The
simplest one, \#1, is the vertical orbit of Section~\ref{sec:vertical} at this
higher energy: computing the vertical orbit directly at $E\approx52.23$ gives
the same period, $T\approx2.861$, and it is linearly stable.
% Checked 2026-11-08 with continue_to_periodic_orbit: T=2.86096, x1*=3.84204,
% max|lambda|=1.00000001 on the full 6-dimensional section.
The others cross the section two to four times per period and have more
complicated shapes. The theorem of \cite{turaev2022} predicts an island of
this kind only around the vertical orbit, and only under the conditions
discussed in Section~\ref{sec:vertical}; the other islands are not covered by
it. We have not measured how much of the regular fraction each island
contributes.

\begin{figure}[H]
  \centering
  \includegraphics[width=\textwidth]{lyapunov_poincare_section_reduced_report.jpg}
  \caption{Poincar\'e section at $E\approx52.23$, each point coloured by the
  Lyapunov exponent $\lambda_{\max}$ of its trajectory (13\,920 trajectories,
  $T=500$). Dark points ($\lambda_{\max}\approx0$) form islands of regular
  motion in the chaotic region. Right: histogram of $\lambda_{\max}$.}
  \label{fig:lyapunov}
\end{figure}

\begin{figure}[H]
  \centering
  \includegraphics[width=\textwidth]{lyapunov_section_periodic_orbits_report.jpg}
  \caption{Stable periodic orbits at the centres of islands, $E\approx52.23$.
  Top: their points on the section. Below: the orbits in the box, with the
  number $k$ of crossings per period and the period $T$.}
  \label{fig:islands}
\end{figure}

\subsection{The vertical orbit: shape and stability as the energy changes}\label{sec:vertical}
This orbit is of the type constructed for a rectangular box in
\cite[Sec.~2.6]{turaev2022}: the two particles move fast in the vertical
direction, each on its own line. Because the motion is centre-symmetric,
particle~1 hits the top wall exactly when particle~2 hits the bottom wall.
This is the ``simultaneous impacts'' case, covered by Theorem~5 of
\cite{turaev2022}. That theorem needs the wall potential to behave like
$(\text{distance to the wall})^{-\beta}$ with $\beta>6$, whereas ours has
$\beta=1$, and it holds only at sufficiently high energy. So the theorem
does not guarantee that the orbit is stable in our system; the results below
are a numerical check of whether it is.

Starting from the vertical orbit at $E\approx402$ (Appendix~\ref{app:cases},
Case~4) with a smaller momentum ($E\approx11.08$), we followed the orbit down
in energy (Figure~\ref{fig:vertical}). Newton converged at every energy. The
results are:
\begin{itemize}
  \item For $E$ from $11.1$ down to $3.88$, all multipliers have $|\lambda|=1$:
        the orbit is linearly stable.
  \item At $E\approx3.32$, $2.77$ and $2.66$ the orbit is unstable, with largest
        multiplier $|\lambda|\approx2.59$, $6.04$ and $5.21$. So it loses
        stability somewhere between $E\approx3.9$ and $3.3$.
  \item At $E\approx2.55$, Newton no longer found the vertical orbit but a
        different, stable orbit in which the particles move along S-shaped
        paths (Figure~\ref{fig:vertical}(d)).
\end{itemize}
We did not study the transition between $E\approx3.9$ and $3.3$ in detail;
this would need much smaller energy steps. So, numerically, the orbit is
stable well outside the conditions of the theorem, down to $E\approx3.9$,
twice the equilibrium energy. Below that it is unstable and has no island of
regular motion around it. The regular fraction rises from $0.20$ at $E=3$ to
$0.25$ at $E=4$ (Figure~\ref{fig:fraction}), the energy range in which the
vertical orbit becomes stable; we have not checked whether this rise comes
from its island.

% Generated by images/make_orbit_figures.py
\begin{figure}[H]
  \centering
  \includegraphics[width=\textwidth]{fig_vertical_family.pdf}
  \caption{One period of the vertical orbit at four energies. (a), (b) stable;
  (c) unstable; (d) at $E\approx2.55$ Newton found a different, S-shaped orbit.
  The black dot is where the orbit crosses the section $y_1=5$.}
  \label{fig:vertical}
\end{figure}

\subsection{The circular orbit: shape and stability as the energy changes}\label{sec:circular}
In this orbit the two particles go round the centre of the box, always opposite
each other. We found it near the trajectory of Case~3 in
Appendix~\ref{app:cases}: Newton started directly from Case~3 did not converge,
but it did converge from a nearby symmetric starting point at the same energy,
using the diagonal $x_1=y_1$ as the section. We then followed this orbit upward
in energy (Figure~\ref{fig:circular}). Newton converged at all 21 energies from
$E\approx2.01$ to $4.02$, and every orbit found is linearly stable. As the
energy increases, the loop grows towards the walls. We also computed the orbit
at much higher energies, $E\approx52$ and $E\approx202$: it has become a
diamond that touches the middle of each wall, and it is still linearly stable.
Since we only computed the orbit at a finite set of energies, we cannot rule
out short energy intervals in between where it is unstable.

At high energy the orbit is thus close to a periodic orbit of a single
particle in a square billiard, the diamond that hits each wall at its
midpoint; the two particles stay far apart on it, so their repulsion is only a
small perturbation. This is the billiard-like motion expected at high energy.
The two particles follow the same curve half a period apart, so this is also a
choreographic motion in the sense of \cite{turaev2022}. Diamond-shaped orbits
are mentioned there as another family of periodic motions in a box on which
such choreographies could be built, but they are not analysed
\cite[Sec.~6]{turaev2022}; we have not checked how our orbit is related to that
family.

\begin{figure}[H]
  \centering
  \includegraphics[width=0.8\textwidth]{fig_circular_family.pdf}
  \caption{One period of the circular orbit at several energies. The two
  particles move on the same curve, half a period apart. (a) For
  $2.0\le E\le4.0$ the loop grows towards the walls. (b) At $E\approx52$ and
  $202$ it is a diamond touching all four walls. All orbits shown are linearly
  stable.}
  \label{fig:circular}
\end{figure}

% ---------------------------------------------------------------------------
\section{Discussion}\label{sec:discussion}

\paragraph{Summary.}
The regular fraction $f(E)$ has three stages. Near the bottom of the wells all
motion is regular, because the wells are nearly harmonic. As the energy
approaches the saddle energy $E_s\approx1.9969$, trajectories can come close to
the saddles and chaos appears, already slightly below $E_s$; above $E_s$, when
particle~1 can move between the wells, chaos takes over almost the whole
energy surface, with a minimum $f\approx0.04$ at $E\approx2.05$. At higher energy $f$ recovers slowly, to about
one third at $E=200$; this regular motion lies on islands around stable
periodic orbits. Of the two orbits we followed, the vertical one is stable for
$E\gtrsim3.9$ and unstable below, and the circular one is stable at every
energy we computed and becomes a billiard-like diamond at high energy.

\paragraph{Relation to Rom-Kedar and Turaev.}
Our results agree with their picture that stable periodic orbits keep the
motion from being ergodic, and extend it in two directions: the vertical orbit
is stable already at moderate energy and with $\beta=1$, outside the
conditions of their theorem; and we measure how much of the energy surface is
regular, which the theorem does not address. Within the energies we computed,
the regular fraction stays well below one.

\paragraph{Limitations.}
All results are for centre-symmetric motions. These form a set of zero volume
in the full phase space, so the regular fraction we measure is a fraction of
the symmetric energy surface, not of the full one, and the islands we see are
islands of the symmetric system. The Lyapunov exponents are computed over a
finite time, so very weak chaos may be counted as regular, and the
classification depends on a threshold; we checked this dependence but cannot
remove it. The stability results for periodic orbits hold at the energies we
computed and say nothing about short intervals in between. Linear stability
also does not prove that an orbit is stable for larger perturbations.

\paragraph{Open questions.}
\begin{itemize}
  \item Does $f(E)$ keep growing above $E=200$, and towards what value?
  \item Are the chaotic trajectories below $E_s$ exactly those that come close
        to the saddles, and how does the width of this chaotic layer depend on
        $E_s-E$?
  \item Which islands make up the regular fraction at high energy, and how
        large is each? In particular, does the rise of $f$ between $E=3$ and
        $4$ come from the island of the vertical orbit?
  \item What happens to the vertical orbit between $E\approx3.9$ and $3.3$,
        and where does the S-shaped orbit come from?
  \item How do the results depend on the exponents $\alpha$, $\beta$, on the
        shape of the box, and on non-symmetric motions?
\end{itemize}

% ---------------------------------------------------------------------------
\appendix

\section{Five example trajectories}\label{app:cases}
Table~\ref{tab:cases} lists five initial conditions, and
Figure~\ref{fig:cases} shows their trajectories and Poincar\'e sections.
\begin{itemize}
  \item Cases~1 and~2 start near the equilibrium and stay near it. Their
        section points lie on smooth curves, so the motion is regular.
  \item In Case~3 the particles go round the centre of the box, always
        opposite each other. The section points form two tiny clusters, which
        suggests a periodic orbit nearby (found in Section~\ref{sec:circular}).
  \item In Case~4 each particle bounces between the top and bottom walls along
        an almost vertical line. This is a periodic orbit: after one period
        $T\approx1.0027$ the state returns to its starting value to within
        $6\times10^{-7}$.
  \item Case~5 starts like Case~4 but at a different $x$. Its section points
        lie on a closed curve around the Case~4 point: the orbit oscillates
        around the stable periodic orbit of Case~4.
\end{itemize}

\begin{table}[H]
\centering
\caption{Initial conditions. Particle~2 always starts at the point reflection
of particle~1 through the centre, with the opposite momentum. $a\approx3.6645$
is the equilibrium position.}
\label{tab:cases}
\begin{tabular}{@{}lllll@{}}
\toprule
Case & $\q_1(0)$ & $\p_1(0)$ & $E$ & Motion\\
\midrule
1 & $(a,a)+5\cdot10^{-4}(1,1)$ & $(10^{-3},0)$ & $1.9877$ & small regular oscillation\\
2 & $(a,a)+0.25\,(1,1)$ & $(0,0.05)$ & $2.0073$ & regular\\
3 & $(a,a)+5\cdot10^{-4}(1,1)$ & $(0.1,-0.1)$ & $2.0077$ & close to a periodic loop\\
4 & $(3.8352,\,5)$ & $(0,20)$ & $402.08$ & periodic, vertical bouncing\\
5 & $(3,\,5)$ & $(0,20)$ & $402.00$ & regular, around Case 4\\
\bottomrule
\end{tabular}
\end{table}

\begin{figure}[H]
  \centering
  \includegraphics[width=\textwidth]{fig_cases.pdf}
  \caption{Top: trajectories of the five cases (particle~1 blue, particle~2
  orange). Bottom: 200 returns to a Poincar\'e section, measured from the
  starting point (red star). The section is taken through the starting point,
  across the largest initial velocity component of particle~1: $x_1=x_1(0)$,
  $p_{x_1}>0$ for Cases~1 and~3, so the axes are $(\Delta y_1,\Delta p_{y_1})$;
  $y_1=y_1(0)$, $p_{y_1}>0$ for Cases~2, 4 and~5, so the axes are
  $(\Delta x_1,\Delta p_{x_1})$. Note the scales: in Case~4 all returns lie
  within $10^{-5}$ of the start.}
  \label{fig:cases}
\end{figure}

\section{Following the circular orbit downwards}\label{app:downwards}
Following the circular orbit to energies \emph{below} the Case~3 energy $E_0$
did not work well. With steps of $0.1\%$ in energy, Newton still converged,
but the period of the orbit it found jumped from $101.8$ to $109.2$, $135.4$
and $140.6$, and at $0.998E_0$ it found an unstable, non-symmetric orbit. So
Newton was jumping between different periodic orbits that lie close together,
instead of following one orbit. At $0.996E_0$ Newton did not converge at all;
with half the energy step it converged again. A continuation method that is
designed to stay on one branch, such as pseudo-arclength continuation
\cite{kuznetsov}, would be the natural next step.

% ---------------------------------------------------------------------------
\section*{Code and reproducibility}
All code is available at
\url{https://github.com/jiajunwu1729/2particle-numerics}. Table~\ref{tab:repro} lists the
script that produces each figure.

\begin{table}[H]
\centering
\caption{Figures and the code that produces them.}
\label{tab:repro}
\begin{tabular}{@{}ll@{}}
\toprule
Figure & Source\\
\midrule
\ref{fig:potential} & \texttt{images/make\_potential\_figure.py}\\
\ref{fig:energy-sections}, \ref{fig:cases} & \texttt{images/make\_section\_figures.py} (same computation as \texttt{report03.ipynb})\\
\ref{fig:threshold}, \ref{fig:fraction} & \texttt{regular\_fraction\_scan.py}, then \texttt{images/make\_threshold\_figure.py}, \texttt{images/make\_fraction\_figure.py}\\
\ref{fig:lyapunov} & \texttt{lyapunov\_poincare\_section.py}, then \texttt{images/compress\_figures.py}\\
\ref{fig:islands} & \texttt{lyapunov\_section\_periodic\_orbits.py}, then \texttt{images/compress\_figures.py}\\
\ref{fig:vertical}, \ref{fig:circular} & \texttt{images/make\_orbit\_figures.py} (periodic points from \texttt{report04.ipynb})\\
\bottomrule
\end{tabular}
\end{table}

\section*{Acknowledgements}
I thank Prof.\ Dmitry Turaev for suggesting this project and for his guidance.

% ---------------------------------------------------------------------------
\begin{thebibliography}{9}
\bibitem{turaev2022}
V.~Rom-Kedar and D.~Turaev,
Stable motions of high energy particles interacting via a repelling potential,
\emph{Comm.\ Math.\ Phys.}\ \textbf{405} (2024), no.~6, 150;
arXiv:2208.14993.
% Verified via Crossref, DOI 10.1007/s00220-024-04946-3 (2026-11-08).

\bibitem{hairer2006}
E.~Hairer, C.~Lubich and G.~Wanner,
\emph{Geometric Numerical Integration}, 2nd ed., Springer, 2006.

\bibitem{lichtenberg1992}
A.~J.~Lichtenberg and M.~A.~Lieberman,
\emph{Regular and Chaotic Dynamics}, 2nd ed., Springer, 1992.

\bibitem{kuznetsov}
Yu.~A.~Kuznetsov, \emph{Elements of Applied Bifurcation Theory}, 3rd ed.,
Springer, 2004.
\end{thebibliography}

\end{document}
